\documentclass[letterpaper]{article}
\usepackage[margin=2cm]{geometry}
\usepackage[english]{babel}
\usepackage{amsfonts}
\usepackage{amsthm}
\usepackage{amsmath}
\usepackage{amssymb}
\usepackage[edges]{forest}
\usepackage{tikz}
\pagestyle{empty}
\usepackage{caption}
\usepackage{enumitem}
\usepackage{setspace}
\usepackage{fdsymbol}
\usepackage{pgfplots}
\usepackage{cancel}
\usepackage{soul}
\usepackage{xcolor}
\usepackage{listings}
\usepackage{graphicx}
\usepackage{parskip}
\usepackage{enumitem}
\setlist{itemsep=-5pt,topsep=-5pt}
\pgfplotsset{compat=1.18} 
 
\begin{document}
\flushleft
\renewcommand{\labelitemiii}{\(\smallblackcircle\)}
\renewcommand{\labelitemii}{\(\smallblackcircle\)}

\lstset{
  basicstyle=\ttfamily,
  mathescape
}

\section*{A5 NP}

\subsection*{Q1.}

\begin{itemize}
  \item R is the set of decidable languages.

  \item P is the set of languages that are decidable in polynomial time.

  \item RE is the set of verifiable languages.

  \item NP is the set of languages that are verifiable in polynomial time.
\end{itemize}

\subsection*{Q2.}

CLIQUE:
\begin{itemize}
  \item Verifier $V$ takes a description of the graph $\langle G \rangle$, $k$,
        and a clique $\langle C \rangle$.
        Checks if the size of $\langle C \rangle$ is at least $k$.
        This can be done in polytime.
        For each pair $u,v$ in $\langle C \rangle$, it checks if there is
        an edge $u,v$ in the graph. This can be done in polytime.
  \item $R = \{ ((G,k), S) : S \text{ is a clique of $G$ with size at least $k$} \}$
\end{itemize}

SUBSET-SUM:
\begin{itemize}
  \item Verifier $V$ takes a list of variables $x_1,...,x_n$,
        $t \in N$, and a subset of the variables $y_1,...,y_m$ such that
        $m \le n$.
        It sums up the variables in the subset and compares it to the total.
        At most $n$ values are summed. This takes polytime.
        Comparison also polytime.
  \item $R = \{ ((x_1,...,x_n,t), (y_1,..., y_m)) : \Sigma_{i=1}^{m} y_i = t \text{, where each $y_i$ is some $x_j$ } \}$
\end{itemize}


QUAD:
\begin{itemize}
  \item Verifier $V$ takes a description of the quadratic polynomials and an assignment.
        For each polynomial, it computes the value using the given assignment and checks if its 0.
        Each polynomial has up to $n^2$ terms, and can be evaluated in polytime.
  \item $R = \{ ((q_1,...,q_n),x): \forall i \in [n], q_i(x)=0  \}$
\end{itemize}

\subsection*{Q3.}

Not in NP because the witness (i.e. $y_i$'s) is not polysized
relative to the input $x$.
Depending on the chosen $f$, witness could be anything.

\subsection*{Q4.}

1. $\text{NTIME}'[T] \subseteq \text{NTIME}[T]$

Given $V'$ that takes a witness with length $T$.
Construct $V$ that takes $x$ and $w$, where $|w| \le T(|x|)$.
If $|w| < T(|x|)$, then pad and run $V'$.
SEE TUTORIAL 6.

2. $\text{NTIME}[T] \subseteq \text{NTIME}'[T'] = \text{NTIME}'[O(T)]$

Given $V$ that takes a witness $w$ of length $\le T$.
Construct $V'$ that takes a witness $w'$ of length $T'$,
where the $w'$ contains $w$ followed by a delimiter $\#$ followed by padding.
$V'$ chops off everything including and following the delimiter and
calls $V$.

WIP

\subsection*{Q5.}

{\setlength{\parskip}{0pt}
  NP is the set of languages that can be verified in polynomial time (deterministic).

  RE is the set of verifiable languages.

  $\Rightarrow NP \subseteq RE$
}


{\setlength{\parskip}{0pt}
  NP is the set of languages that can be decided in polynomial time by a
  non-deterministic machine.

  R is the set of decidable languages.

  $\Rightarrow NP \subseteq R$
}


\subsection*{Q6.}

Let $L \in NP$. Then it has a verifier $V$ that runs in time $O(T(n))$.

Construct D that decides L:
\begin{verbatim}
for each possible witness w st |w| <= T(n):
  run V(x,w)
  if V accepted then accept
  else continue
reject
\end{verbatim}

Total possible number of witnesses is $2^{T(n)}$.
Each call to V is $O(T(n))$.
Total run time is $2^{T(n)} \cdot O(T(n)) = O(2^{T(n)}) = O(2^{n^c})$.
Therefore, L $\in$ EXP.

$\Rightarrow NP \subseteq EXP$

\subsection*{Q7.}
TODO

\subsection*{Q8.}

\begin{enumerate}
  \item no have to prove for every
  \item yes
  \item yes
  \item no that doesnt make sense
\end{enumerate}

For $\ne$, prove 3. For $=$, prove 2.

\subsection*{Q9.}

Same as proof done in class except need to prove every witness
relation for L is polytime verifiable.

See proof for search-to-decision.

Notes:
\begin{itemize}
  \item A language L can have multiple witness relations
  \item A witness relation is a set of $(x,w)$ pairs
  \item An $x$ might have multiple witnesses but the witness relation only needs to contain one of them for each $x \in L$
\end{itemize}

\subsection*{Q10.}

Pick something that's not in NP.

Busy-beaver:
$R = \{ (\langle M \rangle, n) : n \text{ is the max number of steps that M runs for on input } \langle M\rangle \}$

This is not verifiable so $P = NP$ doesn't matter.

\subsection*{Q11.}

$R = \{ (L,x) : L \text{ is a list of 3-constraints satisfied by } x \} $

This is polytime verifiable.

\subsection*{Q12.}

$U = \{ \langle M \rangle , x : M(x) = 1 \}$

TODO

\subsection*{Q13.}

Not necessarily.

Notes:
\begin{itemize}
  \item HALT is not in NP but is NP-hard.
  \item Reduce 3CSP to HALT.
        \begin{itemize}
          \item Given a list of 3-constraints L.
          \item Construct a machine M that iterates over all possible assignments.
                If one of them satisfies L, then accept.
                Otherwise, loop forever.
          \item Can define $\langle M \rangle$ in polynomial time (with respect to the description length of L).
          \item $L \in \text{3CSP} \iff \langle M \rangle \in \text{HALT}$
        \end{itemize}
\end{itemize}

\subsection*{Q14.}

No because every problem in NP can be reduced to an NP-complete one.
Would have to show a reduction of L' to L and that L is in NP.

\subsection*{Q15.}

True? Do something like the CircuitSAT proof. Create a circuit
representing the computation of the TM deciding an arbitrary
problem in P.

\subsection*{Q16.}

P = NP.

If CircuitSAT can be reduced to a problem in P in polynomial time,
then every problem in NP can be reduced to a problem in P in polynomial time.

\subsection*{Q17.}

It goes both ways.

\subsection*{Q18.}

Notes:
\begin{itemize}
  \item Obviously in NP
  \item Reduce 3SAT to 4CSP
  \item Given a 3CNF C of $m$ clauses over $x_1,...,x_n$
  \item Construct a list L of 4-constraints
        \begin{itemize}
          \item For each clause $C_i = l_i \vee l_j \vee l_k$, introduce a new
                variable x and define 2 constraints
                $l_i \vee l_j \vee l_k \vee x = 1$ and $l_i \vee l_j \vee l_k \vee \neg x = 1$.
                The resulting L will consist of variables: $x, x_1,...,x_n$ and have $2m$ constraints.
          \item If C is satisfiable, then for every clause, both corresponding
                constraints are true, because at least one of $l_i, l_j, l_k$ is true since $x$
                and $\neg x$ can't be true at the same time. Therefore, L is satisfiable.
          \item If C is unsatisfiable, there there is some clause where $l_i, l_j, l_k$
                are all false. Consider the corresponding constraints in L, i.e.
                $l_i \vee l_j \vee l_k \vee x = 1$ and $l_i \vee l_j \vee l_k \vee \neg x = 1$
                Since $l_i = l_j = l_k = 0$, if $x=1$ then the second constraint is false,
                if $x=0$ then the first constraint is false. Either way, one of the constraints
                is violated. Therefore, L is unsatisfiable.
        \end{itemize}
\end{itemize}

\subsection*{Q19.}

No. If he did then he proved P=NP which is highly unlikely.

\subsection*{Q20.}

Exponential is fast enough.

\subsection*{Q21.}
TODO

\subsection*{Q22.}

Notes:
\begin{itemize}
  \item Clearly in NP: $V(<C>,x\#y) = 1 \iff C(x) = 1$ and $C(y) = 1$
  \item NP-hard: introduce a dummy variable that doesn't affect the final output
\end{itemize}

WIP

\subsection*{Q23.}

Notes:
\begin{itemize}
  \item Clearly in NP: $V(<C>,x) = 1 \iff C(x) = 1$
  \item NP-hard: Reduce CircuitSAT to Circuit$_{\text{NAND}}$SAT.
        \begin{itemize}
          \item Replace (a AND b) with NOT (a NAND b)
          \item Replace (a OR b) with (a NAND a) NAND (b NAND b)
          \item NOT gates stay the same (or replace (NOT a) with (a NAND a) for fun)
        \end{itemize}
\end{itemize}

WIP

\subsection*{Q24.}

Notes:
\begin{itemize}
  \item The circuit $C_x$ takes input $w \in \{0,1\}^{T}$ where $T=T(n)$.
  \item There are $O(T)$ config layers, one for each step of the computation done by $V(x,w)$
  \item There are a constant number of auxiliary layers between each config
        layer so there are $O(T)$ auxiliary layers
  \item Each config layer had $O(T)$ blocks -- one for each cell on the tape at that config step
  \item Each block has a constant number of nodes so each layer has $O(T)$ nodes.
  \item $O(T)$ layers with $O(T)$ nodes $\implies$ size of the circuit is $O(T^2)$
\end{itemize}

WIP

\subsection*{Q25.}

Notes:
\begin{itemize}
  \item At each configuration layer, only need a constant number of blocks (1?)
        since the location of the head is always known (walks right then left).
        Anywhere that the head is not stays the same, can connect to the gates below.
  \item Since walking right then left, there are around $2T$ steps which is $O(T)$.
  \item Each layer has a constant number of blocks with a constant number of nodes.
  \item $\implies$ size is $O(T)$
\end{itemize}

WIP

\subsection*{Q26.}

All verifiers where the movement of the head is already known.

\subsection*{Q27.}
TODO

\subsection*{Q28.}

Notes:
\begin{itemize}
  \item Clearly in NP: $V(L,x) = 1 \iff $ x is a satisfying assignment for $L$
  \item Show it's NP-hard: WIP
\end{itemize}

\subsection*{Q29.}
Notes:
\begin{itemize}
  \item Clearly in NP: $V( ((S_1,...,S_m), k ), (C_1,...,C_k)) = 1 \iff \cup_{i\in [m]} S_i = \cup_{i\in[k]} C_k$
        \begin{itemize}
          \item For each element in each set $S_i$, check if it's in any of the $C_i$'s
          \item For each element in each set $C_i$, check if it's in any of the $S_i$'s
          \item This can be done in polynomial time with respect to the number of elements (good enough)
        \end{itemize}
  \item Show it's NP-hard: WIP
\end{itemize}

\subsection*{Q30.}
Notes:
\begin{itemize}
  \item Clearly in NP: $V(L,x) = 1 \iff $ x is a satisfying assignment for $L$ where $L$ is a list of 1-in-3 constraints
  \item Show it's NP-hard: WIP
\end{itemize}

\end{document}
